\section{模型评价}

\subsection{模型的优点}

（1）精度与可靠性：我们以有界示向度约束构造定位区域，以区域直径与最小包围圆半径度量其跨度与最坏预测误差。该度量不借助误差的分布假设；最小包围圆同时给出预测位置与误差上界，二者分别是圆心与半径。当半径不超过 20 米时，圆心落入保证清除区域；该上界不随误差分布而变，反映了模型的可靠定位精度。

（2）通用性：我们构建的第二检测点选址模型以接收范围与观测几何作为选点准则，将近距离、方向与无信号三类反馈纳入同一模型。配对测试中，按该模型选点所得预测误差远低于随机选点。三类反馈与不同观测几何共用一套选点准则，选点过程不因反馈类型而更换模型形式。

（3）任务完备性：我们将未知频道巡测、已发现源的补充定位与清除纳入同一任务集合，三者共享同一行动次序并随反馈更新。终止条件同时检查已清除结果与未知频道的剩余区域，任务在两者均无剩余时结束。基于该终止条件，第三、四问的独立测试均完成全部清除并正常退出，清除与排查两类目标均无遗留，完整地完成了既定任务。

（4）可解释性：针对未知发射朝向，我们将接收半径下界与检测点凸包的包围关系结合，给出与源朝向无关的区域排除充分条件。该条件只用检测点位置与接收半径下界两个几何量判定，不含源朝向的估计；检测点调整与路线优化据此始终受完整排查条件约束。判定不依赖统计推断，排除与否可直接由检测点与源的几何关系说明，解释性强。

\subsection{模型的不足}

（1）准则保守性：基于上述最坏情形准则，第二问依据最不利反馈选择检测位置，常见反馈下的平均误差未纳入选点准则。第四问同样以充分条件约束整块区域，要求同一组检测点满足距离与包围关系，该条件会保留部分尚可排除的区域。两部分保守性叠加，巡测工作量因此增加；细化边界区域可缓解这一影响，准则的保守性仍是定位策略的主要代价。

（2）离散近似：第二问采用位置网格加密与方向反馈枚举，第三、四问以有限场景、参数组合及候选路线近似未来任务。这些离散与近似使连续模型的全局最优性无法直接沿用，解的质量仍取决于网格与候选集的规模。

（3）条件依赖：定位模型以源位置固定、检测点坐标准确与示向度误差有界为基础前提。模拟检验采用演练分布与固定平滑误差场，所得精度与耗时只对应这一组测试条件。模型迁移到真实场景时先标定测量误差与接收范围，再检验覆盖情况与任务完成率；所得结果取决于这些前提在真实场景中的成立程度。

